\section{Loading images}
\label{sec:images}

An analysis starts from a folder that holds the image sequence. The \ui{IMAGES} section at the
top of the left sidebar loads the folder, sets the order of the images and lists them. The first
image of the list is the \emph{reference image}: every displacement is measured from it and
reported in its coordinates. On screen it is frame~1; in exported data files and in Python it is
index~0.

\subsection{Opening a folder}

There are three ways to open a folder:
\begin{itemize}
  \item drag the folder onto the drop zone (\ui{Drop image folder or Browse}); if you drop a file
        instead, RAFTcorr opens the folder that contains it;
  \item click the drop zone and choose the folder in the \ui{Select Image Folder} dialog;
  \item choose \menu{File}{Open Image Folder…} (\kbd{Ctrl+I}; \kbd{Cmd+I} on macOS).
\end{itemize}
From a terminal, \file{raftcorr gui FOLDER} starts the application with that folder open
(Section~\ref{sec:cli}).

RAFTcorr reads every image file directly inside the folder. It does not look into subfolders,
and it skips hidden files (names that start with a dot). After loading, the badge next to the
\ui{IMAGES} title shows the number of images, the canvas shows frame~1 fitted to the view, and
the \ui{LOG} in the right sidebar reports the number of images, their size in pixels and the
grey-level scale (next subsection).

\begin{warnbox}
Opening a folder starts a new session. The regions of interest, the results, the probes, the
load data and the model suggestion of the previous images are discarded. Save a project first
if you want to keep them (Section~\ref{sec:projects}).
\end{warnbox}

\subsection{Formats, bit depths and colour}

RAFTcorr reads PNG (\file{.png}), TIFF (\file{.tif}, \file{.tiff}), BMP (\file{.bmp}) and JPEG
(\file{.jpg}, \file{.jpeg}) files; the extension is not case sensitive. Of a multi-page TIFF file
only the first page is read.

The network works on grey levels from 0 to 255. RAFTcorr converts every image with the same
rules, so that the brightness of a speckle does not change from one image to the next because
of the conversion:
\begin{itemize}
  \item \textbf{Colour images} become their luminance ($0.299\,R + 0.587\,G + 0.114\,B$, the
        ITU-R BT.601 weights). An alpha channel is dropped.
  \item \textbf{8-bit images} are used as they are.
  \item \textbf{Images stored with more bits} (16-bit PNG or TIFF files) are scaled by their
        \emph{effective} bit depth: the smallest of 8, 10, 12, 14 and 16 bits that holds the
        brightest pixel of the reference image. A 12-bit camera that writes 16-bit files gives
        values up to 4095, so RAFTcorr maps 4095 (not 65535) to 255 and the speckle keeps its
        contrast.
  \item \textbf{Floating-point images} with values from 0 to 1 are multiplied by 255; any other
        range is divided by the largest value of the reference image and multiplied by 255.
\end{itemize}
The scaled values stay floating-point numbers (they are not rounded back to 8 bits) and are
clipped to the range 0 to 255. The canvas shows the images after this conversion, in grey. The
log line after loading names the scale in use, for example \file{8-bit}, or
\file{12-bit (stored as uint16)} for the camera above.

\begin{notebox}
The scale is decided once, from the reference image, and applied to the whole sequence. If a
later image is brighter than the bit depth of the reference image allows (a 12-bit reference
followed by images with values above 4095), its brightest pixels are clipped to 255. Keep the
exposure constant over the sequence.
\end{notebox}

\subsection{Image size}

All images of a sequence must have the size of the reference image. If one does not, RAFTcorr
shows the \ui{Cannot open the images} dialog and loads nothing. The message names the reference
image with its size and up to five of the images that differ, each with its frame number and its
size, for example \emph{``every frame must have the size of frame 1, a1.png (80 x 64 px);
different: frame 2, a2.png (82 x 64 px)''}.

The models also need a minimum image size: each side must be at least $16\,s$~px, which is
16, 32, 64 and 128~px for RAFTcorr-s1, -s2, -s4 and -s8. A smaller image stops the run with an
error.

\subsection{Sort order}

The order of the list is the order of the sequence, and the first image becomes the reference
image. The \ui{Natural Sort (1, 2, …, 10)} checkbox, on by default, sorts by the numbers inside
the file names: \file{img2.png} comes before \file{img10.png}. Unticked, the list is in plain
alphabetical order, which puts \file{img10.png} before \file{img2.png}. For names with
zero-padded numbers (\file{img\_0002.png}, \file{img\_0010.png}) both orders are the same.

\begin{warnbox}
When changing \ui{Natural Sort} puts the images in a different order, RAFTcorr reloads them and,
like opening a folder, starts a new session: the regions of interest, the results and the probes
are discarded. If there is anything to lose (regions of interest, results, probes, load data or
a running analysis), it asks first with \ui{Change the image order?}; \ui{No} keeps everything
and sets the checkbox back. When the order stays the same, as with zero-padded names, the
checkbox changes nothing else. Set the order before you draw a region of interest.
\end{warnbox}

\subsection{The image list}

\screenshot[0.45\textwidth]{fig-images.png}{The \ui{IMAGES} section after opening a folder of 44
images: the count badge, the drop zone, the \ui{Natural Sort} checkbox and the image list with its
\ui{\#}, \ui{Filename} and \ui{Region} columns.}{fig:images}

The list below the drop zone (Figure~\ref{fig:images}) has three columns:
\begin{description}[leftmargin=1.2cm,style=nextline]
  \item[\#] The frame number, counted from 1 and shown with two digits (\file{01}, \file{02},
    \ldots), the same number as on the frame slider below the canvas.
  \item[Filename] The file name. The rows of the reference frames (the frames that later frames
    are compared with, Section~\ref{sec:workflow}) are tinted with the accent colour and their
    names are bold. Which frames these are depends on the tracking mode, and the highlight
    follows it when you change the mode.
  \item[Region] A button that shows whether the frame has a mask of its own:
    \begin{itemize}
      \item \textcolor{raftred}{\ui{Need}} (red): frame~1 has no region of interest yet. A run
            cannot start without one.
      \item \textcolor{raftgreen}{\ui{Edit}} (green): the frame has a mask of its own.
      \item \ui{Add} (grey): the frame has no mask of its own. This is the normal state of every
            frame after the first; a mask there is optional.
    \end{itemize}
    Clicking the button shows that frame on the canvas and opens its mask for editing
    (Section~\ref{sec:roi-editing}).
\end{description}

Clicking the number or the name of a row shows that frame on the canvas. When the list has the
keyboard focus, \kbd{Up}, \kbd{Down}, \kbd{Page Up}, \kbd{Page Down}, \kbd{Home} and \kbd{End}
move through the frames. \kbd{Ctrl}-click (\kbd{Cmd}-click on macOS) and \kbd{Shift}-click
select several rows.

Right-clicking the selected rows opens a menu:
\begin{description}[leftmargin=1.2cm,style=nextline]
  \item[Import Region of Interest for $n$ frame(s)] Opens a file dialog. Choose exactly as many
    mask files as rows are selected; with any other number nothing is imported and the log says
    that the counts must match. The selected frames, from the first to the last, receive the
    files in the natural order of their names (\file{mask2.png} before \file{mask10.png}),
    whatever order you clicked them in. Check the \ui{Region} buttons afterwards, or use \ui{Batch Import}, which shows every
    pairing before it applies them (Section~\ref{sec:roi}).
  \item[Clear Region of Interest] Removes the masks of the selected frames. The entry tells how
    many of them have one.
  \item[Delete $n$ image(s)] Removes the selected images from the sequence. The files stay on
    disk. The remaining frames are numbered again, their masks follow their images and the
    probes are kept; the results are discarded, because their frame numbers no longer match.
    Deleting frame~1 makes the next image the reference image; the region of interest of the
    deleted frame goes with it, so the new frame~1 keeps only a mask it had of its own.
\end{description}
